
\documentclass{article}
\usepackage{colortbl}
\usepackage{makecell}
\usepackage{multirow}
\usepackage{supertabular}

\begin{document}

\newcounter{utterance}

\centering \large Interaction Transcript for game `function\_detective', experiment `numbers\_passive\_labeled\_pairs', episode 57 with gpt{-}5.4.
\vspace{24pt}

{ \footnotesize  \setcounter{utterance}{1}
\setlength{\tabcolsep}{0pt}
\begin{supertabular}{c@{$\;$}|p{.15\linewidth}@{}p{.15\linewidth}p{.15\linewidth}p{.15\linewidth}p{.15\linewidth}p{.15\linewidth}}
    \# & \multicolumn{2}{c}{Player} && \multicolumn{2}{c}{Game Master} \\
    \hline

    \theutterance \stepcounter{utterance}  
    & & & \multicolumn{4}{p{0.6\linewidth}}{
        \cellcolor[rgb]{0.9,0.9,0.9}{
            \makecell[{{p{\linewidth}}}]{
                \texttt{\tiny{[P1$\langle$GM]}}
                \texttt{I have a hidden Python function for you to figure out.} \\
\\ 
\texttt{Your job is to infer what it does and then submit a correct Python implementation.} \\
\\ 
\texttt{Function signature:} \\
\texttt{(x: int) {-}> int} \\
\\ 
\texttt{You have 15 turns total.} \\
\\ 
\texttt{I reveal candidate input{-}output pairs one at a time and tell you whether each one is correct.} \\
\\ 
\texttt{You must make your best SOLVE attempt no later than the second{-}to{-}last turn.} \\
\texttt{Do not wait until the very end.} \\
\\ 
\texttt{You can use these actions:} \\
\texttt{{-} JSON format described at the end of the prompt} \\
\texttt{{-} SOLVE: ```python ... ```} \\
\\ 
\texttt{Do not add import statements in your guess.} \\
\texttt{Anything you might need is already available, so just write the function body in hypothesis and return the answer.} \\
\\ 
\texttt{Preloaded examples:} \\
\texttt{1. ((99, 5), True)} \\
\\ 
\texttt{Scoring:} \\
\texttt{{-} correctness is required for any non{-}zero score} \\
\texttt{{-} in interactive modes, solving with fewer information requests is better} \\
\\ 
\texttt{Input formatting reminders:} \\
\texttt{{-} strings must be quoted} \\
\texttt{{-} booleans must be True/False} \\
\texttt{{-} use commas between multiple values} \\
\texttt{{-} do not write any extra action tags in the same message} \\
\\ 
\texttt{Format example:} \\
\texttt{You write:} \\
\texttt{\{'current\_mode': <probing/confirming/don't know; choose whichever best describes your current state in the task>, 'input\_rationale': <brief description of why you chose this hypothesis>, 'hypothesis': <either the exact string "unknown" if you have no guess yet OR your current best guess as runnable Python matching the given signature, with a one{-}line docstring stating what it does. Never write placeholder code, if you cannot commit to a guess, use "unknown". Match the given signature.} \\
\texttt{e.g.} \\
\texttt{def f(x):} \\
\texttt{	"""Returns the absolute value of x."""} \\
\texttt{	return x if x >= 0 else {-}x} \\
\\ 
\texttt{Use "unknown" if you have no guess.>, 'confidence\_score': <score between 1{-}100. 100 means highest confidence in your hypothesis>\}} \\
\\ 
\texttt{Output only one valid json per turn.} \\
            }
        }
    }
     \\ \\

    \theutterance \stepcounter{utterance}  
    & \multicolumn{4}{p{0.6\linewidth}}{
        \cellcolor[rgb]{0.561,0.737,0.561}{
            \makecell[{{p{\linewidth}}}]{
                \texttt{\tiny{[P1$\rangle$GM]}}
                \texttt{\{"current\_mode":"probing","input\_rationale":"Probe whether f(100)=1. From the known example f(99)=5, a natural family is digital{-}root{-}like behavior based on the sum of digits; 9+9=18 and 1+8=9, so 5 does not fit plain digital root, but checking a nearby power of 10 helps distinguish digit{-}based rules from divisor/count rules.","hypothesis":"unknown","confidence\_score":12\}} \\
            }
        }
    }
    & & \\ \\

    \theutterance \stepcounter{utterance}  
    & & & \multicolumn{4}{p{0.6\linewidth}}{
        \cellcolor[rgb]{0.9,0.9,0.9}{
            \makecell[{{p{\linewidth}}}]{
                \texttt{\tiny{[P1$\langle$GM]}}
                \texttt{OUTPUT: ((1000, 6), False)} \\
            }
        }
    }
     \\ \\

    \theutterance \stepcounter{utterance}  
    & \multicolumn{4}{p{0.6\linewidth}}{
        \cellcolor[rgb]{0.561,0.737,0.561}{
            \makecell[{{p{\linewidth}}}]{
                \texttt{\tiny{[P1$\rangle$GM]}}
                \texttt{\{"current\_mode":"confirming","input\_rationale":"A strong hypothesis is: return the number of vowels in the English word form of the integer. This fits 99 {-}> 'ninety{-}nine' {-}> 5 vowels, and also explains why 1000 is not 6. Testing 100 is a good discriminator because 'one hundred' has 4 vowels, unlike common arithmetic alternatives.","hypothesis":"def f(x):\textbackslash{}n    \textbackslash{}"\textbackslash{}"\textbackslash{}"Returns the number of vowels in the English word form of x.\textbackslash{}"\textbackslash{}"\textbackslash{}"\textbackslash{}n    n = x\textbackslash{}n    ones = \{\textbackslash{}n        0: 'zero', 1: 'one', 2: 'two', 3: 'three', 4: 'four', 5: 'five', 6: 'six', 7: 'seven', 8: 'eight', 9: 'nine',\textbackslash{}n        10: 'ten', 11: 'eleven', 12: 'twelve', 13: 'thirteen', 14: 'fourteen', 15: 'fifteen', 16: 'sixteen', 17: 'seventeen', 18: 'eighteen', 19: 'nineteen'\textbackslash{}n    \}\textbackslash{}n    tens = \{\textbackslash{}n        20: 'twenty', 30: 'thirty', 40: 'forty', 50: 'fifty', 60: 'sixty', 70: 'seventy', 80: 'eighty', 90: 'ninety'\textbackslash{}n    \}\textbackslash{}n    def words(num):\textbackslash{}n        if num < 20:\textbackslash{}n            return ones{[}num{]}\textbackslash{}n        if num < 100:\textbackslash{}n            if num in tens:\textbackslash{}n                return tens{[}num{]}\textbackslash{}n            return tens{[}num // 10 * 10{]} + '{-}' + ones{[}num \% 10{]}\textbackslash{}n        if num < 1000:\textbackslash{}n            if num \% 100 == 0:\textbackslash{}n                return ones{[}num // 100{]} + ' hundred'\textbackslash{}n            return ones{[}num // 100{]} + ' hundred ' + words(num \% 100)\textbackslash{}n        if num < 1000000:\textbackslash{}n            if num \% 1000 == 0:\textbackslash{}n                return words(num // 1000) + ' thousand'\textbackslash{}n            return words(num // 1000) + ' thousand ' + words(num \% 1000)\textbackslash{}n        if num < 1000000000:\textbackslash{}n            if num \% 1000000 == 0:\textbackslash{}n                return words(num // 1000000) + ' million'\textbackslash{}n            return words(num // 1000000) + ' million ' + words(num \% 1000000)\textbackslash{}n        if num \% 1000000000 == 0:\textbackslash{}n            return words(num // 1000000000) + ' billion'\textbackslash{}n        return words(num // 1000000000) + ' billion ' + words(num \% 1000000000)\textbackslash{}n    s = words(n)\textbackslash{}n    return sum(c in 'aeiou' for c in s)","confidence\_score":61\}} \\
            }
        }
    }
    & & \\ \\

    \theutterance \stepcounter{utterance}  
    & & & \multicolumn{4}{p{0.6\linewidth}}{
        \cellcolor[rgb]{0.9,0.9,0.9}{
            \makecell[{{p{\linewidth}}}]{
                \texttt{\tiny{[P1$\langle$GM]}}
                \texttt{OUTPUT: ((10, 5), True)} \\
            }
        }
    }
     \\ \\

    \theutterance \stepcounter{utterance}  
    & \multicolumn{4}{p{0.6\linewidth}}{
        \cellcolor[rgb]{0.561,0.737,0.561}{
            \makecell[{{p{\linewidth}}}]{
                \texttt{\tiny{[P1$\rangle$GM]}}
                \texttt{\{"current\_mode":"confirming","input\_rationale":"Test the 'number{-}word length chain ending at four' hypothesis. If f(x) returns the final pre{-}4 value in the English letter{-}count chain, then 9 {-}> 'nine' (4 letters), so the answer should be 4.","hypothesis":"def f(x):\textbackslash{}n    \textbackslash{}"\textbackslash{}"\textbackslash{}"Returns the final non{-}4 value in the English letter{-}count chain starting from x.\textbackslash{}"\textbackslash{}"\textbackslash{}"\textbackslash{}n    ones = \{\textbackslash{}n        0: 'zero', 1: 'one', 2: 'two', 3: 'three', 4: 'four', 5: 'five', 6: 'six', 7: 'seven', 8: 'eight', 9: 'nine',\textbackslash{}n        10: 'ten', 11: 'eleven', 12: 'twelve', 13: 'thirteen', 14: 'fourteen', 15: 'fifteen', 16: 'sixteen', 17: 'seventeen', 18: 'eighteen', 19: 'nineteen'\textbackslash{}n    \}\textbackslash{}n    tens = \{\textbackslash{}n        20: 'twenty', 30: 'thirty', 40: 'forty', 50: 'fifty', 60: 'sixty', 70: 'seventy', 80: 'eighty', 90: 'ninety'\textbackslash{}n    \}\textbackslash{}n    def words(n):\textbackslash{}n        if n < 20:\textbackslash{}n            return ones{[}n{]}\textbackslash{}n        if n < 100:\textbackslash{}n            return tens{[}n // 10 * 10{]} if n \% 10 == 0 else tens{[}n // 10 * 10{]} + '{-}' + ones{[}n \% 10{]}\textbackslash{}n        if n < 1000:\textbackslash{}n            return ones{[}n // 100{]} + ' hundred' if n \% 100 == 0 else ones{[}n // 100{]} + ' hundred ' + words(n \% 100)\textbackslash{}n        if n < 1000000:\textbackslash{}n            return words(n // 1000) + ' thousand' if n \% 1000 == 0 else words(n // 1000) + ' thousand ' + words(n \% 1000)\textbackslash{}n        if n < 1000000000:\textbackslash{}n            return words(n // 1000000) + ' million' if n \% 1000000 == 0 else words(n // 1000000) + ' million ' + words(n \% 1000000)\textbackslash{}n        return words(n // 1000000000) + ' billion' if n \% 1000000000 == 0 else words(n // 1000000000) + ' billion ' + words(n \% 1000000000)\textbackslash{}n    def letter\_count(n):\textbackslash{}n        return sum(c.isalpha() for c in words(n))\textbackslash{}n    n = letter\_count(x)\textbackslash{}n    while letter\_count(n) != 4:\textbackslash{}n        n = letter\_count(n)\textbackslash{}n    return n","confidence\_score":68\}} \\
            }
        }
    }
    & & \\ \\

    \theutterance \stepcounter{utterance}  
    & & & \multicolumn{4}{p{0.6\linewidth}}{
        \cellcolor[rgb]{0.9,0.9,0.9}{
            \makecell[{{p{\linewidth}}}]{
                \texttt{\tiny{[P1$\langle$GM]}}
                \texttt{OUTPUT: (({-}2, 6), False)} \\
            }
        }
    }
     \\ \\

    \theutterance \stepcounter{utterance}  
    & \multicolumn{4}{p{0.6\linewidth}}{
        \cellcolor[rgb]{0.561,0.737,0.561}{
            \makecell[{{p{\linewidth}}}]{
                \texttt{\tiny{[P1$\rangle$GM]}}
                \texttt{\{"current\_mode":"confirming","input\_rationale":"Check the 'last letter{-}count before reaching 4 in the English number{-}word chain' hypothesis on a one{-}step case. For 9, the word is 'nine' with 4 letters, so the result should be 4 if this is correct.","hypothesis":"def f(x):\textbackslash{}n    \textbackslash{}"\textbackslash{}"\textbackslash{}"Returns the last value in the English letter{-}count chain before it reaches 4.\textbackslash{}"\textbackslash{}"\textbackslash{}"\textbackslash{}n    neg = x < 0\textbackslash{}n    n = {-}x if neg else x\textbackslash{}n    ones = \{\textbackslash{}n        0: 'zero', 1: 'one', 2: 'two', 3: 'three', 4: 'four', 5: 'five', 6: 'six', 7: 'seven', 8: 'eight', 9: 'nine',\textbackslash{}n        10: 'ten', 11: 'eleven', 12: 'twelve', 13: 'thirteen', 14: 'fourteen', 15: 'fifteen', 16: 'sixteen', 17: 'seventeen', 18: 'eighteen', 19: 'nineteen'\textbackslash{}n    \}\textbackslash{}n    tens = \{\textbackslash{}n        20: 'twenty', 30: 'thirty', 40: 'forty', 50: 'fifty', 60: 'sixty', 70: 'seventy', 80: 'eighty', 90: 'ninety'\textbackslash{}n    \}\textbackslash{}n    def words(m):\textbackslash{}n        if m < 20:\textbackslash{}n            return ones{[}m{]}\textbackslash{}n        if m < 100:\textbackslash{}n            return tens{[}m // 10 * 10{]} if m \% 10 == 0 else tens{[}m // 10 * 10{]} + '{-}' + ones{[}m \% 10{]}\textbackslash{}n        if m < 1000:\textbackslash{}n            return ones{[}m // 100{]} + ' hundred' if m \% 100 == 0 else ones{[}m // 100{]} + ' hundred ' + words(m \% 100)\textbackslash{}n        if m < 1000000:\textbackslash{}n            return words(m // 1000) + ' thousand' if m \% 1000 == 0 else words(m // 1000) + ' thousand ' + words(m \% 1000)\textbackslash{}n        if m < 1000000000:\textbackslash{}n            return words(m // 1000000) + ' million' if m \% 1000000 == 0 else words(m // 1000000) + ' million ' + words(m \% 1000000)\textbackslash{}n        return words(m // 1000000000) + ' billion' if m \% 1000000000 == 0 else words(m // 1000000000) + ' billion ' + words(m \% 1000000000)\textbackslash{}n    def letter\_count(m):\textbackslash{}n        s = ('negative ' if neg and m == n else '') + words(m)\textbackslash{}n        return sum(c.isalpha() for c in s)\textbackslash{}n    c = letter\_count(n)\textbackslash{}n    prev = c\textbackslash{}n    while c != 4:\textbackslash{}n        prev = c\textbackslash{}n        c = letter\_count(c)\textbackslash{}n    return prev","confidence\_score":73\}} \\
            }
        }
    }
    & & \\ \\

    \theutterance \stepcounter{utterance}  
    & & & \multicolumn{4}{p{0.6\linewidth}}{
        \cellcolor[rgb]{0.9,0.9,0.9}{
            \makecell[{{p{\linewidth}}}]{
                \texttt{\tiny{[P1$\langle$GM]}}
                \texttt{OUTPUT: (({-}1, 5), True)} \\
            }
        }
    }
     \\ \\

    \theutterance \stepcounter{utterance}  
    & \multicolumn{4}{p{0.6\linewidth}}{
        \cellcolor[rgb]{0.561,0.737,0.561}{
            \makecell[{{p{\linewidth}}}]{
                \texttt{\tiny{[P1$\rangle$GM]}}
                \texttt{\{"current\_mode":"confirming","input\_rationale":"Differentiate 'return the last value before 4 in the English letter{-}count chain' from variants that always return 5. If x=4, its word has 4 letters, so the correct output should be 4 under the current hypothesis.","hypothesis":"def f(x):\textbackslash{}n    \textbackslash{}"\textbackslash{}"\textbackslash{}"Returns the last value in the English letter{-}count chain before it reaches 4.\textbackslash{}"\textbackslash{}"\textbackslash{}"\textbackslash{}n    orig\_neg = x < 0\textbackslash{}n    if x < 0:\textbackslash{}n        x = {-}x\textbackslash{}n    ones = \{\textbackslash{}n        0: 'zero', 1: 'one', 2: 'two', 3: 'three', 4: 'four', 5: 'five', 6: 'six', 7: 'seven', 8: 'eight', 9: 'nine',\textbackslash{}n        10: 'ten', 11: 'eleven', 12: 'twelve', 13: 'thirteen', 14: 'fourteen', 15: 'fifteen', 16: 'sixteen', 17: 'seventeen', 18: 'eighteen', 19: 'nineteen'\textbackslash{}n    \}\textbackslash{}n    tens = \{\textbackslash{}n        20: 'twenty', 30: 'thirty', 40: 'forty', 50: 'fifty', 60: 'sixty', 70: 'seventy', 80: 'eighty', 90: 'ninety'\textbackslash{}n    \}\textbackslash{}n    def words(n):\textbackslash{}n        if n < 20:\textbackslash{}n            return ones{[}n{]}\textbackslash{}n        if n < 100:\textbackslash{}n            return tens{[}n // 10 * 10{]} if n \% 10 == 0 else tens{[}n // 10 * 10{]} + '{-}' + ones{[}n \% 10{]}\textbackslash{}n        if n < 1000:\textbackslash{}n            return ones{[}n // 100{]} + ' hundred' if n \% 100 == 0 else ones{[}n // 100{]} + ' hundred ' + words(n \% 100)\textbackslash{}n        if n < 1000000:\textbackslash{}n            return words(n // 1000) + ' thousand' if n \% 1000 == 0 else words(n // 1000) + ' thousand ' + words(n \% 1000)\textbackslash{}n        if n < 1000000000:\textbackslash{}n            return words(n // 1000000) + ' million' if n \% 1000000 == 0 else words(n // 1000000) + ' million ' + words(n \% 1000000)\textbackslash{}n        return words(n // 1000000000) + ' billion' if n \% 1000000000 == 0 else words(n // 1000000000) + ' billion ' + words(n \% 1000000000)\textbackslash{}n    def count\_letters\_of\_original():\textbackslash{}n        s = words(x)\textbackslash{}n        if orig\_neg:\textbackslash{}n            s = 'negative ' + s\textbackslash{}n        return sum(c.isalpha() for c in s)\textbackslash{}n    def count\_letters(n):\textbackslash{}n        return sum(c.isalpha() for c in words(n))\textbackslash{}n    c = count\_letters\_of\_original()\textbackslash{}n    prev = c\textbackslash{}n    while c != 4:\textbackslash{}n        prev = c\textbackslash{}n        c = count\_letters(c)\textbackslash{}n    return prev","confidence\_score":86\}} \\
            }
        }
    }
    & & \\ \\

    \theutterance \stepcounter{utterance}  
    & & & \multicolumn{4}{p{0.6\linewidth}}{
        \cellcolor[rgb]{0.9,0.9,0.9}{
            \makecell[{{p{\linewidth}}}]{
                \texttt{\tiny{[P1$\langle$GM]}}
                \texttt{OUTPUT: (({-}1000, 6), False)} \\
            }
        }
    }
     \\ \\

    \theutterance \stepcounter{utterance}  
    & \multicolumn{4}{p{0.6\linewidth}}{
        \cellcolor[rgb]{0.561,0.737,0.561}{
            \makecell[{{p{\linewidth}}}]{
                \texttt{\tiny{[P1$\rangle$GM]}}
                \texttt{\{"current\_mode":"confirming","input\_rationale":"The current hypothesis predicts a special case at x=4: 'four' already has 4 letters, so the chain starts at 4 and the function should return 4 rather than 5. This cleanly tests whether the function is the pre{-}4 value in the number{-}letter{-}count chain.","hypothesis":"def f(x):\textbackslash{}n    \textbackslash{}"\textbackslash{}"\textbackslash{}"Returns the last value in the English number{-}letter{-}count chain before reaching 4.\textbackslash{}"\textbackslash{}"\textbackslash{}"\textbackslash{}n    neg = x < 0\textbackslash{}n    x = {-}x if x < 0 else x\textbackslash{}n    ones = \{\textbackslash{}n        0: 'zero', 1: 'one', 2: 'two', 3: 'three', 4: 'four', 5: 'five', 6: 'six', 7: 'seven', 8: 'eight', 9: 'nine',\textbackslash{}n        10: 'ten', 11: 'eleven', 12: 'twelve', 13: 'thirteen', 14: 'fourteen', 15: 'fifteen', 16: 'sixteen', 17: 'seventeen', 18: 'eighteen', 19: 'nineteen'\textbackslash{}n    \}\textbackslash{}n    tens = \{20: 'twenty', 30: 'thirty', 40: 'forty', 50: 'fifty', 60: 'sixty', 70: 'seventy', 80: 'eighty', 90: 'ninety'\}\textbackslash{}n    def words(n):\textbackslash{}n        if n < 20:\textbackslash{}n            return ones{[}n{]}\textbackslash{}n        if n < 100:\textbackslash{}n            return tens{[}n // 10 * 10{]} if n \% 10 == 0 else tens{[}n // 10 * 10{]} + '{-}' + ones{[}n \% 10{]}\textbackslash{}n        if n < 1000:\textbackslash{}n            return ones{[}n // 100{]} + ' hundred' if n \% 100 == 0 else ones{[}n // 100{]} + ' hundred ' + words(n \% 100)\textbackslash{}n        if n < 1000000:\textbackslash{}n            return words(n // 1000) + ' thousand' if n \% 1000 == 0 else words(n // 1000) + ' thousand ' + words(n \% 1000)\textbackslash{}n        if n < 1000000000:\textbackslash{}n            return words(n // 1000000) + ' million' if n \% 1000000 == 0 else words(n // 1000000) + ' million ' + words(n \% 1000000)\textbackslash{}n        return words(n // 1000000000) + ' billion' if n \% 1000000000 == 0 else words(n // 1000000000) + ' billion ' + words(n \% 1000000000)\textbackslash{}n    def count\_letters\_in\_original():\textbackslash{}n        s = words(x)\textbackslash{}n        if neg:\textbackslash{}n            s = 'negative ' + s\textbackslash{}n        return sum(c.isalpha() for c in s)\textbackslash{}n    def count\_letters(n):\textbackslash{}n        return sum(c.isalpha() for c in words(n))\textbackslash{}n    c = count\_letters\_in\_original()\textbackslash{}n    if c == 4:\textbackslash{}n        return 4\textbackslash{}n    prev = c\textbackslash{}n    while c != 4:\textbackslash{}n        prev = c\textbackslash{}n        c = count\_letters(c)\textbackslash{}n    return prev","confidence\_score":89\}} \\
            }
        }
    }
    & & \\ \\

    \theutterance \stepcounter{utterance}  
    & & & \multicolumn{4}{p{0.6\linewidth}}{
        \cellcolor[rgb]{0.9,0.9,0.9}{
            \makecell[{{p{\linewidth}}}]{
                \texttt{\tiny{[P1$\langle$GM]}}
                \texttt{OUTPUT: ((0, 5), True)} \\
            }
        }
    }
     \\ \\

    \theutterance \stepcounter{utterance}  
    & \multicolumn{4}{p{0.6\linewidth}}{
        \cellcolor[rgb]{0.561,0.737,0.561}{
            \makecell[{{p{\linewidth}}}]{
                \texttt{\tiny{[P1$\rangle$GM]}}
                \texttt{\{"current\_mode":"confirming","input\_rationale":"Distinguish a constant{-}5 function from the English letter{-}count{-}chain hypothesis. For x=17, 'seventeen' has 9 letters and 'nine' has 4, so that hypothesis predicts 9 rather than 5.","hypothesis":"def f(x):\textbackslash{}n    \textbackslash{}"\textbackslash{}"\textbackslash{}"Returns 5.\textbackslash{}"\textbackslash{}"\textbackslash{}"\textbackslash{}n    return 5","confidence\_score":44\}} \\
            }
        }
    }
    & & \\ \\

    \theutterance \stepcounter{utterance}  
    & & & \multicolumn{4}{p{0.6\linewidth}}{
        \cellcolor[rgb]{0.9,0.9,0.9}{
            \makecell[{{p{\linewidth}}}]{
                \texttt{\tiny{[P1$\langle$GM]}}
                \texttt{OUTPUT: (({-}99, 6), False)} \\
            }
        }
    }
     \\ \\

    \theutterance \stepcounter{utterance}  
    & \multicolumn{4}{p{0.6\linewidth}}{
        \cellcolor[rgb]{0.561,0.737,0.561}{
            \makecell[{{p{\linewidth}}}]{
                \texttt{\tiny{[P1$\rangle$GM]}}
                \texttt{\{"current\_mode":"confirming","input\_rationale":"Test a case where the English letter{-}count{-}chain hypothesis is not 5: for x=17, 'seventeen' has 9 letters, and 9 {-}> 'nine' has 4, so this hypothesis predicts f(17)=9. Candidate pair: ((17, 9), True).","hypothesis":"def f(x):\textbackslash{}n    \textbackslash{}"\textbackslash{}"\textbackslash{}"Returns the last value in the English number{-}letter{-}count chain before reaching 4.\textbackslash{}"\textbackslash{}"\textbackslash{}"\textbackslash{}n    neg = x < 0\textbackslash{}n    x = {-}x if x < 0 else x\textbackslash{}n    ones = \{\textbackslash{}n        0: 'zero', 1: 'one', 2: 'two', 3: 'three', 4: 'four', 5: 'five', 6: 'six', 7: 'seven', 8: 'eight', 9: 'nine',\textbackslash{}n        10: 'ten', 11: 'eleven', 12: 'twelve', 13: 'thirteen', 14: 'fourteen', 15: 'fifteen', 16: 'sixteen', 17: 'seventeen', 18: 'eighteen', 19: 'nineteen'\textbackslash{}n    \}\textbackslash{}n    tens = \{20: 'twenty', 30: 'thirty', 40: 'forty', 50: 'fifty', 60: 'sixty', 70: 'seventy', 80: 'eighty', 90: 'ninety'\}\textbackslash{}n    def words(n):\textbackslash{}n        if n < 20:\textbackslash{}n            return ones{[}n{]}\textbackslash{}n        if n < 100:\textbackslash{}n            return tens{[}n // 10 * 10{]} if n \% 10 == 0 else tens{[}n // 10 * 10{]} + '{-}' + ones{[}n \% 10{]}\textbackslash{}n        if n < 1000:\textbackslash{}n            return ones{[}n // 100{]} + ' hundred' if n \% 100 == 0 else ones{[}n // 100{]} + ' hundred ' + words(n \% 100)\textbackslash{}n        if n < 1000000:\textbackslash{}n            return words(n // 1000) + ' thousand' if n \% 1000 == 0 else words(n // 1000) + ' thousand ' + words(n \% 1000)\textbackslash{}n        if n < 1000000000:\textbackslash{}n            return words(n // 1000000) + ' million' if n \% 1000000 == 0 else words(n // 1000000) + ' million ' + words(n \% 1000000)\textbackslash{}n        return words(n // 1000000000) + ' billion' if n \% 1000000000 == 0 else words(n // 1000000000) + ' billion ' + words(n \% 1000000000)\textbackslash{}n    def count\_letters\_original():\textbackslash{}n        s = words(x)\textbackslash{}n        if neg:\textbackslash{}n            s = 'negative ' + s\textbackslash{}n        return sum(c.isalpha() for c in s)\textbackslash{}n    def count\_letters(n):\textbackslash{}n        return sum(c.isalpha() for c in words(n))\textbackslash{}n    c = count\_letters\_original()\textbackslash{}n    if c == 4:\textbackslash{}n        return 4\textbackslash{}n    prev = c\textbackslash{}n    while c != 4:\textbackslash{}n        prev = c\textbackslash{}n        c = count\_letters(c)\textbackslash{}n    return prev","confidence\_score":84\}} \\
            }
        }
    }
    & & \\ \\

    \theutterance \stepcounter{utterance}  
    & & & \multicolumn{4}{p{0.6\linewidth}}{
        \cellcolor[rgb]{0.9,0.9,0.9}{
            \makecell[{{p{\linewidth}}}]{
                \texttt{\tiny{[P1$\langle$GM]}}
                \texttt{OUTPUT: ((1, 4), True)} \\
            }
        }
    }
     \\ \\

    \theutterance \stepcounter{utterance}  
    & \multicolumn{4}{p{0.6\linewidth}}{
        \cellcolor[rgb]{0.561,0.737,0.561}{
            \makecell[{{p{\linewidth}}}]{
                \texttt{\tiny{[P1$\rangle$GM]}}
                \texttt{\{"current\_mode":"confirming","input\_rationale":"New hypothesis: f(x) counts the length of the 'four is magic' chain starting from the character length of the decimal representation of x (including a leading minus sign). This predicts {-}99 has representation length 3, giving chain 3 {-}> 5 {-}> 4, so the total length including the original x is 4.","hypothesis":"def f(x):\textbackslash{}n    \textbackslash{}"\textbackslash{}"\textbackslash{}"Returns the length of the chain from x to 4 by first taking len(str(x)), then repeatedly taking the letter count of the English name of the current number.\textbackslash{}"\textbackslash{}"\textbackslash{}"\textbackslash{}n    ones = \{\textbackslash{}n        0: 'zero', 1: 'one', 2: 'two', 3: 'three', 4: 'four', 5: 'five', 6: 'six', 7: 'seven', 8: 'eight', 9: 'nine',\textbackslash{}n        10: 'ten', 11: 'eleven', 12: 'twelve', 13: 'thirteen', 14: 'fourteen', 15: 'fifteen', 16: 'sixteen', 17: 'seventeen', 18: 'eighteen', 19: 'nineteen'\textbackslash{}n    \}\textbackslash{}n    tens = \{20: 'twenty', 30: 'thirty', 40: 'forty', 50: 'fifty', 60: 'sixty', 70: 'seventy', 80: 'eighty', 90: 'ninety'\}\textbackslash{}n    def words(n):\textbackslash{}n        if n < 20:\textbackslash{}n            return ones{[}n{]}\textbackslash{}n        if n < 100:\textbackslash{}n            return tens{[}n // 10 * 10{]} if n \% 10 == 0 else tens{[}n // 10 * 10{]} + '{-}' + ones{[}n \% 10{]}\textbackslash{}n        if n < 1000:\textbackslash{}n            return ones{[}n // 100{]} + ' hundred' if n \% 100 == 0 else ones{[}n // 100{]} + ' hundred ' + words(n \% 100)\textbackslash{}n        if n < 1000000:\textbackslash{}n            return words(n // 1000) + ' thousand' if n \% 1000 == 0 else words(n // 1000) + ' thousand ' + words(n \% 1000)\textbackslash{}n        if n < 1000000000:\textbackslash{}n            return words(n // 1000000) + ' million' if n \% 1000000 == 0 else words(n // 1000000) + ' million ' + words(n \% 1000000)\textbackslash{}n        if n < 1000000000000:\textbackslash{}n            return words(n // 1000000000) + ' billion' if n \% 1000000000 == 0 else words(n // 1000000000) + ' billion ' + words(n \% 1000000000)\textbackslash{}n        return str(n)\textbackslash{}n    def letter\_count(n):\textbackslash{}n        s = words(n)\textbackslash{}n        return sum(c.isalpha() for c in s)\textbackslash{}n    count = 1\textbackslash{}n    n = len(str(x))\textbackslash{}n    while True:\textbackslash{}n        count += 1\textbackslash{}n        if n == 4:\textbackslash{}n            return count\textbackslash{}n        n = letter\_count(n)","confidence\_score":91\}} \\
            }
        }
    }
    & & \\ \\

    \theutterance \stepcounter{utterance}  
    & & & \multicolumn{4}{p{0.6\linewidth}}{
        \cellcolor[rgb]{0.9,0.9,0.9}{
            \makecell[{{p{\linewidth}}}]{
                \texttt{\tiny{[P1$\langle$GM]}}
                \texttt{OUTPUT: (({-}10, 6), False)} \\
            }
        }
    }
     \\ \\

    \theutterance \stepcounter{utterance}  
    & \multicolumn{4}{p{0.6\linewidth}}{
        \cellcolor[rgb]{0.561,0.737,0.561}{
            \makecell[{{p{\linewidth}}}]{
                \texttt{\tiny{[P1$\rangle$GM]}}
                \texttt{\{"current\_mode":"confirming","input\_rationale":"Test whether 1 is uniquely special among small positive integers. If f(2)=5, that strongly supports a simple predicate{-}like rule where only x=1 maps to 4 and most other values map to 5.","hypothesis":"def f(x):\textbackslash{}n    \textbackslash{}"\textbackslash{}"\textbackslash{}"Returns 4 for 1 and 5 otherwise.\textbackslash{}"\textbackslash{}"\textbackslash{}"\textbackslash{}n    return 4 if x == 1 else 5","confidence\_score":58\}} \\
            }
        }
    }
    & & \\ \\

    \theutterance \stepcounter{utterance}  
    & & & \multicolumn{4}{p{0.6\linewidth}}{
        \cellcolor[rgb]{0.9,0.9,0.9}{
            \makecell[{{p{\linewidth}}}]{
                \texttt{\tiny{[P1$\langle$GM]}}
                \texttt{OUTPUT: ((2, 3), True)} \\
            }
        }
    }
     \\ \\

    \theutterance \stepcounter{utterance}  
    & \multicolumn{4}{p{0.6\linewidth}}{
        \cellcolor[rgb]{0.561,0.737,0.561}{
            \makecell[{{p{\linewidth}}}]{
                \texttt{\tiny{[P1$\rangle$GM]}}
                \texttt{\{"current\_mode":"probing","input\_rationale":"I want to test whether the outputs for small nonnegative integers follow a descending simple pattern from 0,1,2. Checking x=3 is highly informative: under several candidate chain/word{-}length rules it could be 5, 4, or 2, so it should sharply narrow the hypothesis space.","hypothesis":"unknown","confidence\_score":19\}} \\
            }
        }
    }
    & & \\ \\

    \theutterance \stepcounter{utterance}  
    & & & \multicolumn{4}{p{0.6\linewidth}}{
        \cellcolor[rgb]{0.9,0.9,0.9}{
            \makecell[{{p{\linewidth}}}]{
                \texttt{\tiny{[P1$\langle$GM]}}
                \texttt{OUTPUT: (({-}50, {-}1), False)} \\
            }
        }
    }
     \\ \\

    \theutterance \stepcounter{utterance}  
    & \multicolumn{4}{p{0.6\linewidth}}{
        \cellcolor[rgb]{0.561,0.737,0.561}{
            \makecell[{{p{\linewidth}}}]{
                \texttt{\tiny{[P1$\rangle$GM]}}
                \texttt{\{"current\_mode":"confirming","input\_rationale":"Current best simple hypothesis consistent with all revealed data is that only the small positives 1 and 2 are exceptional: f(1)=4, f(2)=3, and otherwise f(x)=5. A value at x=3 or x=4 would sharply confirm or refute this.","hypothesis":"def f(x):\textbackslash{}n    \textbackslash{}"\textbackslash{}"\textbackslash{}"Returns 4 for 1, 3 for 2, and 5 for all other integers.\textbackslash{}"\textbackslash{}"\textbackslash{}"\textbackslash{}n    if x == 1:\textbackslash{}n        return 4\textbackslash{}n    if x == 2:\textbackslash{}n        return 3\textbackslash{}n    return 5","confidence\_score":41\}} \\
            }
        }
    }
    & & \\ \\

    \theutterance \stepcounter{utterance}  
    & & & \multicolumn{4}{p{0.6\linewidth}}{
        \cellcolor[rgb]{0.9,0.9,0.9}{
            \makecell[{{p{\linewidth}}}]{
                \texttt{\tiny{[P1$\langle$GM]}}
                \texttt{OUTPUT: ((11, 5), True)} \\
            }
        }
    }
     \\ \\

    \theutterance \stepcounter{utterance}  
    & \multicolumn{4}{p{0.6\linewidth}}{
        \cellcolor[rgb]{0.561,0.737,0.561}{
            \makecell[{{p{\linewidth}}}]{
                \texttt{\tiny{[P1$\rangle$GM]}}
                \texttt{\{"current\_mode":"confirming","input\_rationale":"A minimal hypothesis still consistent with all confirmed data is that 1 and 2 are the only exceptions, with all other tested inputs mapping to 5. Testing x=3 is the sharpest confirmation of that simple rule.","hypothesis":"def f(x):\textbackslash{}n    \textbackslash{}"\textbackslash{}"\textbackslash{}"Returns 4 for 1, 3 for 2, and 5 for all other integers.\textbackslash{}"\textbackslash{}"\textbackslash{}"\textbackslash{}n    if x == 1:\textbackslash{}n        return 4\textbackslash{}n    if x == 2:\textbackslash{}n        return 3\textbackslash{}n    return 5","confidence\_score":46\}} \\
            }
        }
    }
    & & \\ \\

    \theutterance \stepcounter{utterance}  
    & & & \multicolumn{4}{p{0.6\linewidth}}{
        \cellcolor[rgb]{0.9,0.9,0.9}{
            \makecell[{{p{\linewidth}}}]{
                \texttt{\tiny{[P1$\langle$GM]}}
                \texttt{OUTPUT: (({-}30, 12), False)} \\
            }
        }
    }
     \\ \\

    \theutterance \stepcounter{utterance}  
    & \multicolumn{4}{p{0.6\linewidth}}{
        \cellcolor[rgb]{0.561,0.737,0.561}{
            \makecell[{{p{\linewidth}}}]{
                \texttt{\tiny{[P1$\rangle$GM]}}
                \texttt{\{"current\_mode":"probing","input\_rationale":"The confirmed data are still consistent with two very different possibilities: a genuinely structured rule or a mostly{-}constant function with only tiny exceptions. Testing x=3 is the cleanest discriminator, since many plausible rules predict 2, 4, or 5 there.","hypothesis":"def f(x):\textbackslash{}n    \textbackslash{}"\textbackslash{}"\textbackslash{}"Returns 4 for 1, 3 for 2, and 5 for all other integers.\textbackslash{}"\textbackslash{}"\textbackslash{}"\textbackslash{}n    if x == 1:\textbackslash{}n        return 4\textbackslash{}n    if x == 2:\textbackslash{}n        return 3\textbackslash{}n    return 5","confidence\_score":38\}} \\
            }
        }
    }
    & & \\ \\

    \theutterance \stepcounter{utterance}  
    & & & \multicolumn{4}{p{0.6\linewidth}}{
        \cellcolor[rgb]{0.9,0.9,0.9}{
            \makecell[{{p{\linewidth}}}]{
                \texttt{\tiny{[P1$\langle$GM]}}
                \texttt{OUTPUT: ((23, 4), True)} \\
            }
        }
    }
     \\ \\

    \theutterance \stepcounter{utterance}  
    & \multicolumn{4}{p{0.6\linewidth}}{
        \cellcolor[rgb]{0.561,0.737,0.561}{
            \makecell[{{p{\linewidth}}}]{
                \texttt{\tiny{[P1$\rangle$GM]}}
                \texttt{SOLVE: ```python} \\
\texttt{def f(x):} \\
\texttt{    """Return 5 for negative integers; otherwise use the first digit and digit count."""} \\
\texttt{    if x < 0:} \\
\texttt{        return 5} \\
\texttt{    s = str(x)} \\
\texttt{    return abs(int(s{[}0{]}) {-} 5) + len(s) {-} 1} \\
\texttt{```} \\
            }
        }
    }
    & & \\ \\

    \theutterance \stepcounter{utterance}  
    & & & \multicolumn{2}{p{0.3\linewidth}}{
        \cellcolor[rgb]{0.9,0.9,0.9}{
            \makecell[{{p{\linewidth}}}]{
                \texttt{\tiny{[GM$|$GM]}}
                \texttt{Game Verdict: LOSS} \\
            }
        }
    }
    & & \\ \\

    \theutterance \stepcounter{utterance}  
    & & & \multicolumn{2}{p{0.3\linewidth}}{
        \cellcolor[rgb]{0.9,0.9,0.9}{
            \makecell[{{p{\linewidth}}}]{
                \texttt{\tiny{[GM$|$GM]}}
                \texttt{Incorrect.} \\
\texttt{Failed Cases:} \\
\texttt{Input: {[}1000{]} | Expected: 5 | Got: 7} \\
\texttt{Input: {[}{-}2{]} | Expected: 4 | Got: 5} \\
\texttt{Input: {[}{-}1000{]} | Expected: 4 | Got: 5} \\
            }
        }
    }
    & & \\ \\

    \theutterance \stepcounter{utterance}  
    & & & \multicolumn{2}{p{0.3\linewidth}}{
        \cellcolor[rgb]{0.9,0.9,0.9}{
            \makecell[{{p{\linewidth}}}]{
                \texttt{\tiny{[GM$|$GM]}}
                \texttt{def absolute\_modulus\_gap\_eleven(x: int) {-}> int:} \\
\texttt{    return abs((x \% 11) {-} 5)} \\
            }
        }
    }
    & & \\ \\

    \theutterance \stepcounter{utterance}  
    & & & \multicolumn{2}{p{0.3\linewidth}}{
        \cellcolor[rgb]{0.9,0.9,0.9}{
            \makecell[{{p{\linewidth}}}]{
                \texttt{\tiny{[GM$|$GM]}}
                \texttt{def f(x):} \\
\texttt{    """Return 5 for negative integers; otherwise use the first digit and digit count."""} \\
\texttt{    if x < 0:} \\
\texttt{        return 5} \\
\texttt{    s = str(x)} \\
\texttt{    return abs(int(s{[}0{]}) {-} 5) + len(s) {-} 1} \\
            }
        }
    }
    & & \\ \\

    \theutterance \stepcounter{utterance}  
    & & & \multicolumn{2}{p{0.3\linewidth}}{
        \cellcolor[rgb]{0.9,0.9,0.9}{
            \makecell[{{p{\linewidth}}}]{
                \texttt{\tiny{[GM$|$GM]}}
                \texttt{Mode: passive\_labeled\_pairs} \\
\texttt{Accuracy: 0.250} \\
\texttt{Information requests: 14/15} \\
\texttt{Revealed passive examples: 15} \\
\texttt{Efficiency: 0.133} \\
\texttt{Slack: 1} \\
\texttt{Observed{-}pair consistency: 1.000} \\
\texttt{Membership balance: 8 true / 7 false} \\
\texttt{Parse errors: 0} \\
\texttt{Runtime errors: 0} \\
            }
        }
    }
    & & \\ \\

\end{supertabular}
}

\end{document}
